% ID: FIS-OND-COR-001
% Titolo: Corde di chitarra e frequenza del suono
% Disciplina: Fisica
% Area: Onde
% Argomento: Onde su corde
% Sottoargomento: Frequenza e caratteristiche della corda
% Classe: Quarta liceo scientifico
% Difficolta: 2
% Tipo: quesito_teorico
% Risultato: soluzione da aggiungere
% Tempo_stimato: 6 min
% Competenze: interpretazione fisica, collegamento tra modello e fenomeno reale, argomentazione
% Prerequisiti: frequenza, onde stazionarie, velocita di propagazione su una corda
% Tag: fisica, onde, onde_su_corde, chitarra, frequenza
% Fonte: verifica_fisica_onde_anonimizzata
% Autore: Riccardo Carli
% Licenza: CC-BY-SA-4.0

\begin{esercizio}
Le corde di una chitarra che producono le note piu gravi hanno un diametro
maggiore rispetto a quelle che producono suoni piu acuti.

Spiega qual e il motivo fisico che sta alla base di questa differenza.
\end{esercizio}

\begin{soluzione}
% Soluzione da aggiungere.
\end{soluzione}

